% ---------------------------------------------------------------------------
% The Lund string. (a) The colour field between a quark and an antiquark is
% a tube of constant energy per length, kappa = 1 GeV/fm, so V(r) = kappa r.
% (b) The world sheet of a two-parton string in its rest frame: the two ends
% fly apart at the speed of light, the string between them breaks at
% vertices that lie near a hyperbola of constant proper time, and every
% piece between two neighbouring vertices is one hadron. (c) The numbers of
% our event (the -v log, "the colour-singlet strings of the record").
% ---------------------------------------------------------------------------
\begin{tikzpicture}[x=1cm,y=1cm,
  every node/.style={font=\footnotesize},
  q/.style={circle,inner sep=0pt,minimum size=2.4mm,fill=red!80!black},
  aq/.style={circle,inner sep=0pt,minimum size=2.4mm,fill=blue!80!black},
  vtx/.style={circle,inner sep=0pt,minimum size=1.8mm,fill=black},
  lab/.style={font=\scriptsize,align=center},
  head/.style={font=\small\bfseries}]
% ---------------- (a) the flux tube ----------------------------------------------------------
\begin{scope}
\node[head,anchor=west] at (-0.3,3.3) {(a) the rope: $V(r)=\kappa r$};
\fill[orange!35] (0.6,1.75) rectangle (4.4,2.25);
\foreach \x in {0.8,1.2,...,4.2} {\draw[orange!70!black,thin] (\x,1.75) -- (\x+0.2,2.25);}
\node[q] (q) at (0.5,2.0) {}; \node[aq] (aq) at (4.5,2.0) {};
\node[lab,above=1pt of q] {$q$}; \node[lab,above=1pt of aq] {$\bar q$};
\draw[-{Stealth}] (q) -- ++(-0.6,0); \draw[-{Stealth}] (aq) -- ++(0.6,0);
\draw[{Stealth}-{Stealth},thin] (0.6,1.4) -- (4.4,1.4) node[midway,below,lab] {$r$};
\node[lab,anchor=north,text width=5.4cm] at (2.5,0.95) {energy stored: $\kappa r$ with $\kappa=1\GeV/\mathrm{fm}=0.197\GeV^2$; a $1\,\mathrm{fm}$ string holds $1\GeV$, enough for a $u\bar u$ pair.\\[2pt] \textcolor{black!65}{our $28\GeV$ two-parton string would stretch to $14+14\,\mathrm{fm}$ before turning round; it breaks long before}};
\end{scope}
% ---------------- (b) the world sheet --------------------------------------------------------
\begin{scope}[xshift=9.2cm,yshift=-0.3cm]
\node[head,anchor=west] at (-2.9,3.6) {(b) the world sheet breaks};
\begin{scope}[x=0.32cm,y=0.32cm]
\fill[orange!18] (0,0) -- (-8,8) -- (8,8) -- cycle;
\draw[-{Stealth}] (-8.5,0) -- (8.8,0) node[below left,lab] {$x$ [fm]};
\draw[-{Stealth}] (0,0) -- (0,9.2) node[left,lab] {$t$ [fm]};
\draw[red!80!black,thick] (0,0) -- (-8,8); \draw[blue!80!black,thick] (0,0) -- (8,8);
\node[lab,text=red!80!black] at (-7.3,8.7) {$q$}; \node[lab,text=blue!80!black] at (7.3,8.7) {$\bar q$};
% hyperbola of constant proper time tau = 1.5 fm: t = sqrt(x^2 + 2.25)
\draw[black!60,dashed] plot[domain=-7.5:7.5,samples=61] (\x,{sqrt(\x*\x+2.25)});
\node[lab,text=black!65,anchor=west] at (5.0,2.0) {$\Gamma=\kappa^2\tau^2=x^+x^-$};
% vertices at rapidities -2.0 -1.0 -0.2 0.6 1.4 2.2
\foreach \px/\pt in {-5.44/5.64,-1.76/2.31,-0.30/1.53,0.96/1.78,2.86/3.22,6.72/6.88} {
  \node[vtx] at (\px,\pt) {};
  \draw[black!70] (\px,\pt) -- ++(1.1,1.1); \draw[black!70] (\px,\pt) -- ++(-1.1,1.1);}
\node[lab,anchor=south] at (-3.6,6.2) {hadron};
\node[lab,anchor=south] at (-1.0,3.2) {hadron};
\node[lab,anchor=south] at (0.3,2.9) {$\cdots$};
\node[lab,anchor=south] at (4.8,7.4) {hadron};
\node[lab,anchor=north west,text width=5.6cm,text=black!65] at (-8.6,-1.3) {each break: a $q'\bar q'$ pair from the vacuum; the piece between two neighbouring vertices is one hadron, the end pieces carry the original $q$ and $\bar q$};
\end{scope}
\end{scope}
% ---------------- (c) the numbers --------------------------------------------------------------
\begin{scope}[xshift=12.9cm]
\node[head,anchor=west] at (-0.3,3.3) {(c) our event};
\node[lab,anchor=north west,text width=5.2cm,align=left] at (-0.3,3.0) {
211 partons, every colour tag once as colour and once as anticolour\\[2pt]
$\to$ 33 colour singlets: 29 strings and 4 ministrings, no closed gluon loop; 48 partons joined to a neighbour (pair mass within $0.3\GeV$ of the endpoint masses)\\[2pt]
total string mass $3071\GeV$; the heaviest, $1036\GeV$, runs from a remnant $d$ to a remnant $\bar u$; the $b$ quark's, $69.3\GeV$, has four gluons\\[2pt]
$\to$ 365 primary hadrons; a two-parton string gives $n\approx0.1+3.0\ln(W/\GeV)$: 7 from $28\GeV$, 18 from $1036\GeV$};
\end{scope}
\end{tikzpicture}
